\section{Three-prime support and mixed-full Campana points}
\label{sec:three-prime-campana}

This section records two complementary descriptions of the first locus not
covered by the two-prime theorem.  The first is an exact support
decomposition.  The second keeps multiplicities rather than support size and
isolates the log-general-type Campana locus studied most recently by
Browning--Verzobio~\cite{BrowningVerzobio2026}.

\subsection{The support-three gate}

Let $a+b=c$ be a positive primitive triple with
$\omega(abc)\leq 3$.  Pairwise coprimality gives
\[
 \omega(abc)=\omega(a)+\omega(b)+\omega(c).
\]
If $a,b>1$, every term on the right is positive.  Hence equality holds in
the support bound, and there are distinct primes $p,q,r$ and positive
integers $x,y,z$ such that
\begin{equation}\label{eq:three-prime-fermat}
                     a=p^x,\qquad b=q^y,\qquad c=r^z.
\end{equation}
If one input is one, the support split between the two consecutive nonunits
is exactly $(1,1)$, $(1,2)$, or $(2,1)$.  Thus, after exchanging the inputs,
the remaining equations are
\begin{equation}\label{eq:three-prime-unit-branches}
 1+p^x=q^y,\qquad 1+p^x=q^yr^z,\qquad
 1+p^xq^y=r^z,
\end{equation}
with distinct displayed prime bases.  This proves an exhaustion rather than
a density reduction: the two semiprime branches in
\eqref{eq:three-prime-unit-branches} remain moving Diophantine problems.

The all-nonunit branch has the following useful rigidity.

\begin{theorem}[Common-exponent rigidity]\label{thm:three-prime-common-exp}
Suppose that distinct primes satisfy
$p^x+q^y=r^z$ and $d=\gcd(x,y,z)>1$.  Then, up to exchanging the inputs,
\[
                         2^4+3^2=5^2.
\]
\end{theorem}

\begin{proof}
Set $X=p^{x/d}$, $Y=q^{y/d}$, and $Z=r^{z/d}$.  Then
$X^d+Y^d=Z^d$.  Fermat's Last Theorem~\cite{Wiles1995,TaylorWiles1995}
forces $d=2$.  In the resulting primitive Pythagorean triple the even leg is
a power of two, say $X=2^A$.  Euclid's parametrization gives
\[
 X=2mn,\qquad Y=m^2-n^2,\qquad Z=m^2+n^2,
\]
where $m,n$ are coprime and of opposite parity.  Since $mn$ is a power of
two, one has $n=1$ and $m=2^k$.  Now
\[
 Y=(2^k-1)(2^k+1)
\]
is a prime power, while the two odd factors are coprime.  The smaller factor
must be one, so $k=1$ and $(X,Y,Z)=(4,3,5)$.  Squaring proves the assertion.
\end{proof}

For $x,y,z\geq2$, sorting the exponents shows that
$1/x+1/y+1/z\geq1$ occurs only for
\[
 (2,2,n),\ (2,3,3),\ (2,3,4),\ (2,3,5),\
 (3,3,3),\ (2,4,4),\ (2,3,6).
\]
Elementary difference-of-powers arguments exclude the full signatures
$(2,3,3)$ and $(2,3,6)$ in every output orientation.  Theorem
\ref{thm:three-prime-common-exp} excludes $(3,3,3)$ and $(2,4,4)$.
For $(2,2,n)$, the odd-exponent orientation reduces to
$S^2+4=R^n$; the theorem of Arif--Abu Muriefah
\cite{ArifAbuMuriefah2001} then yields the second of exactly two solutions
\[
                 2^4+3^2=5^2,\qquad 2^2+11^2=5^3.
\]
The other orientations follow by factoring a difference of squares; even
$n$ follows from Theorem~\ref{thm:three-prime-common-exp}.

Consequently, if $c>\rad(abc)$ in an all-nonunit support-three triple, then
the exponent gcd is one and the all-nonlinear remainder consists only of
$(2,2,3)$, $(2,3,4)$, $(2,3,5)$, or signatures with reciprocal sum below
one.  The first case is the genuine triple
$2^2+11^2=5^3$, for which $125>110$.  Linear exponents are essential:
\[
                 7+181^2=2^{15},\qquad
 \frac{2^{15}}{2\cdot7\cdot181}>12.
\]
This example refutes a putative extension $c\leq\rad(abc)$ to support
three; it is not an abc counterexample.  The unresolved branches are the two
semiprime unit equations, linear exponents, the prime-base signatures
$(2,3,4)$ and $(2,3,5)$, and the union over varying hyperbolic signatures.
Darmon--Granville~\cite{DarmonGranville1995} gives finiteness for each fixed
hyperbolic signature, not uniformly over that union.

\subsection{Different fullness exponents}

An integer $n>0$ is $m$-full when every prime divisor occurs with valuation
at least $m$.  Let $P=(a,b,c)$ be a primitive abc point and suppose that
$a,b,c$ are respectively $p$-, $q$-, and $r$-full, where $p,q,r>0$.  Put
\[
       \sigma_{p,q,r}=\frac1p+\frac1q+\frac1r.
\]

\begin{theorem}[Mixed-full radical compression]
\label{thm:mixed-full-compression}
One has
\begin{align}
 \rad(abc)^{pqr}&\le c^{qr+pr+pq},\label{eq:mixed-full-natural}\\
 \log\rad(abc)&\le \sigma_{p,q,r}\log c.\label{eq:mixed-full-log}
\end{align}
Moreover,
\[
 \sigma_{p,q,r}<1\quad\Longleftrightarrow\quad
 qr+pr+pq<pqr.
\]
\end{theorem}

\begin{proof}
Fullness gives $\rad(a)^p\mid a$, and similarly for $b,c$.
Raise these three divisibilities to $qr,pr,pq$, multiply them, and use
radical multiplicativity across a primitive triple.  Since $a,b<c$, this
gives
\[
 \rad(abc)^{pqr}
 \le a^{qr}b^{pr}c^{pq}
 \le c^{qr+pr+pq},
\]
which proves \eqref{eq:mixed-full-natural}.  Taking the three logarithmic
inequalities separately gives
\[
 \log\rad(abc)
 \le \frac{\log a}{p}+\frac{\log b}{q}+\frac{\log c}{r}
 \le \sigma_{p,q,r}\log c,
\]
proving \eqref{eq:mixed-full-log}.  Clearing the positive denominator $pqr$
proves the final equivalence.
\end{proof}

\begin{corollary}[A strict Campana disproof gate]
\label{cor:mixed-full-disproof-gate}
Fix $p,q,r>0$ with $\sigma_{p,q,r}<1$.  The abc conjecture implies a
uniform height bound for every primitive abc point with the three prescribed
fullness conditions.  Conversely, an unbounded family of such points
implies the negation of the standard abc conjecture.
\end{corollary}

\begin{proof}
Choose $\eps>0$ so that $(1+\eps)\sigma_{p,q,r}<1$.  If abc holds with
constant $C_\eps$, Theorem~\ref{thm:mixed-full-compression} gives
\[
 \bigl(1-(1+\eps)\sigma_{p,q,r}\bigr)\log c\le C_\eps.
\]
The coefficient is positive, so $c$ is bounded.  This proves the first
claim.  The same displayed inequality contradicts unbounded height and
proves the second.
\end{proof}

Browning--Verzobio~\cite[Theorems 1.1 and 1.2]{BrowningVerzobio2026}
study exactly these points.  If $p=r+u$, $q=r+v$, with fixed $u\geq v\geq0$,
they prove, for all sufficiently large $r$,
\[
 N(B)\ll B^{1/p+1/q-\eta_{u,v}(r)+\eps},\qquad
 \eta_{u,v}(r)=r^{-2}+O_{u,v}(r^{-5/2})>0.
\]
They also obtain a power saving for every $r\geq2$ in the special family
$(p,q,r)=(r+2,r+1,r)$.  This is a genuine improvement over the trivial
$B^{1/p+1/q}$ count and applies precisely to the low-radical locus in
Corollary~\ref{cor:mixed-full-disproof-gate}.  Its exponent remains positive,
so it proves neither the predicted boundedness nor the existence of an
unbounded family.  The determinant-method and function-field inputs of that
paper are cited mathematical results and are not declared as Lean axioms.
